%% Section 6: Discussion
\section{Discussion, Implications, and Limitations}
\label{sec:discussion}

\subsection{Implications for Researchers}

The formal trade-off framework (\S\ref{sec:tradeoff}) reveals that
the research community has disproportionately optimised for
$\mathcal{D}_1$ (semantic precision) while under-investing in
$\mathcal{D}_4$ (update cost) and $\mathcal{D}_5$ (scalability).
This imbalance is partly an artefact of benchmark
design: SWE-bench~\cite{jimenez2024swebench} and
RepoBench~\cite{liu2024repobench} are static snapshots that do not
penalise stale representations.  HumanEvo~\cite{zheng2025humanevo}
is a notable exception, and its adoption as a primary benchmark
for future PAR-3 evaluations is strongly recommended.

\subsection{Implications for Practitioners}

System designers should treat the five-dimension framework as a
\emph{requirements elicitation} tool: before selecting a paradigm,
explicitly specify hard constraints on each dimension (\textit{e.g.}\ 
``latency $\leq$ 200 ms'', ``update cost $\leq$ 5 s'',
``repository size up to 50K files'').  The decision guidelines
G1--G5 (\S\ref{subsec:tradeoff:guidelines}) then yield a unique
paradigm recommendation or a hybrid strategy.

\subsection{Threats to Validity}

\textit{Internal validity}: the 1--5 dimension scores are based on
evidence from primary studies that use heterogeneous experimental
setups.  We report evidence directly from each study rather than
averaging incompatible numbers, but a meta-analytic synthesis would
provide stronger quantitative ground.

\textit{External validity}: the corpus skews toward Python
repositories (SWE-bench is Python-only) and may not generalise to
statically-typed languages (Java, C++) where PAR-3 tooling is more
mature and PAR-2 baselines are weaker.

\textit{Construct validity}: our five dimensions are not
independent; token efficiency ($\mathcal{D}_2$) and construction
latency ($\mathcal{D}_3$) are correlated.  Future work should
apply factor analysis to validate the dimension structure.
